% Part: many-valued-logic
% Chapter: sequent-calculus
% Section: introduction

\documentclass[../../../include/open-logic-section]{subfiles}

\begin{document}

\olfileid[fr]{mvl}{seq}{int}

\olsection{Introduction}

Le calcul des séquents de la logique classique est un système de
!!{derivation} simple et efficace. Lorsqu'une logique plurivalente est
définie par une matrice dont l'ensemble des valeurs de vérité~$V$ est
fini, on peut également lui associer un calcul des séquents. Pour voir
comment le construire, partons du sens des séquents et de la forme des
règles d'inférence dans le cas classique.

Rappelons qu'un séquent
\begin{align*}
    !A_1, \dots, !A_m & \Sequent !B_1, \dots, !B_n
\intertext{peut s'interpréter comme la !!{formula}}
    (!A_1 \land \cdots \land !A_m) & \lif (!B_1 \lor \cdots \lor
!B_n)
\end{align*}
La conjonction vide vaut $\True$ et la disjonction vide vaut $\False$.
Autrement dit, !!^a{valuation}~$\pAssign{v}$ \emph{satisfait} un
séquent $\Gamma \Sequent \Delta$ si et seulement si
$\pValue{v}(!A)=\False$ pour au moins une formule $!A\in\Gamma$,
ou $\pValue{v}(!A)=\True$ pour au moins une formule $!A\in\Delta$.
Dans cette interprétation, les séquents initiaux
$!A \Sequent !A$ sont toujours satisfaits : on a soit
$\pValue{v}(!A)=\True$, soit $\pValue{v}(!A)=\False$.

Voici les règles d'inférence pour l'implication dans~$\Log{LK}$,
en omettant les formules de contexte $\Gamma$ et~$\Delta$ :

\begin{defish}
    \Axiom$ \fCenter !A$
    \Axiom$ !B \fCenter $
    \RightLabel{\LeftR{\lif}}
    \BinaryInf$ !A \lif !B \fCenter $
    \DisplayProof
    \hfill
    \Axiom$ !A \fCenter !B$
    \RightLabel{\RightR{\lif}}
    \UnaryInf$ \fCenter !A \lif !B $
    \DisplayProof
\end{defish}

Selon l'interprétation sémantique précédente, la règle
$\LeftR{\lif}$ dit que, si $\pValue{v}(!A)=\True$ et
$\pValue{v}(!B)=\False$, alors
$\pValue{v}(!A\lif !B)=\False$. De même, la règle
$\RightR{\lif}$ dit que, si $\pValue{v}(!A)=\False$ ou
$\pValue{v}(!B)=\True$, alors
$\pValue{v}(!A\lif !B)=\True$. En fait, ces deux implications
sémantiques sont des équivalences. Pour les formulations usuelles des
règles $\LeftR{\land}$ et $\RightR{\lor}$, les implications
correspondantes ne seraient pas des équivalences. Mais elles le
deviennent avec les variantes suivantes :

\begin{defish}
    \Axiom$!A, !B, \Gamma \fCenter \Delta$
    \RightLabel{\LeftR{\land}}
    \UnaryInf$!A \land !B, \Gamma \fCenter \Delta$
    \DisplayProof
    \hfill
    \Axiom$ \Gamma \fCenter \Delta, !A, !B$
    \RightLabel{\RightR{\lor}}
    \UnaryInf$ \Gamma \fCenter \Delta, !A \lor !B$
    \DisplayProof
\end{defish}

En appliquant cette idée à une logique à $n$ valeurs, on obtient un
calcul des séquents à $n$ places plutôt qu'à deux, une pour chaque
valeur de vérité. Pour une logique à trois valeurs avec
$V=\{\False,\Undef,\True\}$, un séquent est une expression
$\Gamma\mid\Pi\mid\Delta$. Il est satisfait par
!!a{valuation}~$\pAssign{v}$ si et seulement si
$\pValue{v}(!A)=\False$ pour au moins une formule $!A\in\Gamma$,
ou $\pValue{v}(!A)=\True$ pour au moins une formule $!A\in\Delta$,
ou encore $\pValue{v}(!A)=\Undef$ pour au moins une formule
$!A\in\Pi$. Dans un séquent initial $!A\mid !A\mid !A$, quelle que
soit la valeur de~$!A$, la place qui lui correspond contient~$!A$.
Toute valuation satisfait donc ce séquent.

\end{document}
